% CANONICAL SECTION ORDER (do not reorder):
%   PROFILE -> WORK EXPERIENCE -> PROJECTS -> EDUCATION -> SKILLS
% SKILLS MUST be the last section, not second. See SKILL.md Step 6 and apply.md
% Step 7 for the workflow that depends on this order.
%
% Location heading in the document body: city only, no parentheticals
% (e.g. "Glasgow, UK", not "Glasgow, UK (willing to relocate; UK Graduate
% visa, no sponsorship required)"). See CLAUDE.md "No relocation/visa notes
% in CV headings" for the full rule.
\documentclass[11pt,a4paper, sans-serif]{article}
\renewcommand{\familydefault}{\sfdefault}
\usepackage[margin=0.75in]{geometry}
% colorlinks=true renders email, LinkedIn, GitHub, and project URLs in blue.
% Do NOT switch to [hidelinks] — that hides all links and breaks the
% highlighted-hyperlink contract that the workflow depends on.
\usepackage[colorlinks=true,
            linkcolor=blue,
            urlcolor=blue,
            citecolor=blue]{hyperref}
% \usepackage{fontawesome}
\usepackage{needspace}
\usepackage{titlesec}
\usepackage{enumitem}
\usepackage{xcolor}


\titleformat{\section}{\large\scshape\bfseries}{}{0em}{}[\titlerule]
\titlespacing{\section}{0pt}{10pt}{5pt}

\begin{document}

% Header
\begin{center}
    {\LARGE \bfseries Firstname Lastname} \\[5pt]
    Glasgow, UK $|$ +44 0000000000 $|$ \href{mailto:candidate@example.com}{candidate@example.com} $|$ 
    \href{https://linkedin.com}{LinkedIn} $|$ \href{https://github.com}{GitHub}
\end{center}

\vspace{10pt}

% Profile
\section*{PROFILE}
\noindent AI Engineer specialising in Generative AI, Large Language Models, Retrieval-Augmented Generation (RAG), and production AI systems. Experienced building and deploying cloud-native AI applications using Python, FastAPI, LangChain, Docker, AWS, LangSmith, and Ragas, taking solutions from research through production. Delivered measurable improvements including 24\% higher retrieval recall, 48\% lower LLM token usage, and 97\% faster large-scale data processing, while applying responsible AI practices through bias auditing and continuous LLM evaluation.

% Work Experience
\section*{WORK EXPERIENCE}

\noindent\textbf{NeuraSearch Laboratory} $|$ \textbf{AI Research Intern} $|$ Glasgow, UK (Remote) \hfill Oct 2025 - Feb 2026
\begin{itemize}[leftmargin=*, nosep]
    \item Built an LLM-based retrieval pipeline with OpenAI and LangChain, from prompt design and document ingestion through retrieval tuning and output assessment.
    \item Compared multiple language models and retrieval strategies using LangSmith and Ragas across large evaluation datasets, identifying the best-performing configuration and improving retrieval recall by 24\% while maintaining answer faithfulness.
    \item Cut token overhead by 48\% with RECOMP compression while maintaining measured faithfulness scores.
    \item Automated document processing with a reusable Python pipeline, reducing processing time per 100-document batch by 97\%.
    \item Documented model behaviour, experiment configurations, and findings so workflows stayed reproducible after handover.
\end{itemize}


\noindent\textbf{CL Techno} $|$ \textbf{(Volunteer) AI Engineer} $|$ London, UK (Remote) \hfill Feb 2024 - Feb 2025

\begin{itemize}[leftmargin=*, nosep]
    \item Built an automated machine learning workflow using Python, PyTorch and Docker, connecting preprocessing, training, evaluation and inference into a single end-to-end pipeline that eliminated manual intervention. 
    \item Improved a production computer vision model through iterative tuning and deployment optimisation, increasing inference accuracy from 84\% to 86\% while supporting live client deployments. 
    \item Collaborated fully remotely through written communication and GitHub, independently delivering documented code, technical updates and reusable solutions across distributed teams.
\end{itemize}


\noindent\textbf{Gamaka AI} $|$ \textbf{Data Analyst Intern} $|$ Pune, India \hfill Jul 2023 - Dec 2023
\begin{itemize}[leftmargin=*, nosep]
    \item Developed a customer spend-prediction model for e-commerce users using regression techniques, achieving 98\% validation accuracy.
    \item Analysed business data using Python, SQL, and statistical techniques, presenting actionable insights to non-technical stakeholders through Power BI dashboards.
\end{itemize}

% Projects
\section*{PROJECTS}
\noindent\textbf{Violence Detection using Deep Learning} \\
\noindent Developed and deployed a Vision Transformer (ViT) image classification system using PyTorch, achieving 98.82\% accuracy on approximately 3,000 images. Benchmarked multiple deep learning architectures and optimised model performance through transfer learning and systematic evaluation across Precision, Recall, F1-score and ROC-AUC.

\noindent\textbf{Archival Retrieval Analytics System (RAG)} \\
\noindent Built a large-scale data analytics pipeline processing 121k+ documents using Python-based preprocessing and machine learning evaluation methods. Applied statistical benchmarking techniques to compare retrieval approaches and measure model performance improvements.

% Education
\section*{EDUCATION}

\noindent\textbf{MSc Advanced Computer Science with AI} $|$ University of Strathclyde, UK \hfill Sep 2024 - Sep 2025 
\textbf{Dissertation:} Auditing Bias in Retrieval-Augmented Generation Systems Using Cultural Heritage Archives. \\
\textbf{Modules:} Machine Learning, Deep Learning, Big Data Tools and Techniques, Quantitative Methods for AI.

% Skills
\section*{SKILLS}
\begin{itemize}[leftmargin=*, nosep]
    \item \textbf{Gen AI \& LLM Engineering:} LangChain, RAG, OpenAI API, prompt engineering, MCP protocol, agentic workflows, LLM fine-tuning.
    \item \textbf{MLOps \& LLMOps:} LangSmith observability, Ragas evaluation (faithfulness, relevancy, hallucination rate), experiment tracking, monitoring, evaluation.
    \item \textbf{Deep Learning:} PyTorch, TensorFlow 2.x, Transformers, transfer learning, fine-tuning.
    \item \textbf{Cloud \& Tools:} AWS (S3, EC2, SageMaker), Docker, Linux, GitHub.
    \item \textbf{Responsible AI \& AI-Assisted Development:} bias auditing in RAG systems, fairness-aware evaluation, agentic development with Claude Code and Cursor.
\end{itemize}

\end{document}